\documentclass[12pt, a4paper]{article}
\usepackage{amsmath}
\usepackage{amssymb}
\usepackage{geometry}
\usepackage{hyperref}
\usepackage{booktabs}
\usepackage{array}
\geometry{margin=1in}

\title{\textbf{Free-Space Six-Degree-of-Freedom Magnetic Levitation\\for Two-Needle Hand Knitting Mechanics}}
\date{}

\begin{document}

\maketitle

\begin{abstract}
Existing magnetically driven knitting systems constrain needles mechanically
within a machine bed, with magnetic actuation controlling only vertical
displacement. This paper proposes a fundamentally different approach: a
two-needle system modeled on hand knitting mechanics, where full six-degree-of-freedom
(6-DOF) needle actuation is achieved purely through magnetic means, without
mechanical guides or onboard controllers. The needle design employs an aluminum
core for eddy current torque generation and a non-continuous magnetic strip for
discrete pole levitation control. PWM-sequenced solenoid arrays manage all
translational and rotational degrees of freedom. Camera-based AI monitoring
provides yarn feedback in place of contact sensors.
\end{abstract}

%------------------------------------------------------------
\section{Introduction}
%------------------------------------------------------------

Automated knitting has long been dominated by industrial machines that use
mechanical cam systems to drive arrays of needles. These systems are effective
but mechanically complex, noisy, and prone to needle breakage due to physical
contact between components.

Recent work has explored replacing mechanical cam drives with magnetic actuation.
Zhang et al.\ \cite{zhang2019} proposed a non-contact magnetic array needle driving
method where permanent magnet needles are driven vertically by an electromagnet
array. Subsequent work \cite{zuo2021, xiong2023, zhang2025} developed this approach
further, optimizing needle gauge, magnetic shielding between needles, and control
strategies including PID and sliding mode controllers. However, all of these systems
share a fundamental architectural constraint: needles are mechanically guided within
a machine bed, and magnetic actuation controls only the vertical (axial) displacement
of each needle. The architecture remains that of an industrial weft knitting machine
with a magnetic upgrade.

This paper proposes a departure from that architecture. Rather than scaling up an
industrial machine, the goal is to replicate the mechanics of hand knitting --- two
needles, free-space actuation, full 6-DOF control --- using a solenoid array and
purpose-designed needles. The key contributions of this work are:

\begin{itemize}
    \item A two-needle free-space magnetic actuation architecture modeled on hand
    knitting mechanics.
    \item A purpose-designed needle where each component serves a specific
    degree of freedom.
    \item A PWM sequencing strategy that simultaneously manages torque generation
    and cross-interference suppression.
    \item A camera-based AI monitoring approach for yarn feedback without
    contact sensors.
\end{itemize}

%------------------------------------------------------------
\section{System Overview}
%------------------------------------------------------------

The system consists of two identical actuator units, each comprising a solenoid
array and a knitting needle. The two units are positioned facing each other,
separated by an insulating boundary, in a configuration that mirrors the
two-handed mechanics of hand knitting.

Each solenoid array is a $p \times q$ grid of individually addressable miniature
solenoids with silicon-steel laminated cores. Each knitting needle sits above its
respective array and is actuated purely through magnetic forces and torques
generated by the array. No mechanical guides constrain the needle's motion.

The knitted fabric is collected in a wooden tray positioned below the work area.
The tray height is set according to the breaking stress of the needle and the
linear mass density of the knitting pattern, ensuring the needle is never
overloaded by the weight of accumulated fabric.

%------------------------------------------------------------
\section{Needle Design}
%------------------------------------------------------------

The needle is designed so that each component directly serves one or more
degrees of freedom. No element is arbitrary.

\subsection{Geometry}
The needle has a cylindrical shaft. The leading end is conical, optimized
for yarn manipulation during loop formation. The trailing end is flat,
which maximizes the surface area available for eddy current generation
and therefore torque transfer.

\subsection{Aluminum Core}
The core of the needle is aluminum. Aluminum is non-ferromagnetic, meaning
it does not retain or redirect magnetic fields passively. In regions of the
needle where no magnetic strip is present, the aluminum core is magnetically
inert to static or slowly varying fields. This eliminates the need for
additional passive shielding on the shaft.

Under a time-varying magnetic field, eddy currents are induced in the
aluminum core. These eddy currents interact with the driving magnetic field
to produce torque, which controls the rotational degrees of freedom
(XY, YZ, ZX planes). The flat trailing end maximizes this effect by
increasing the effective cross-sectional area for eddy current circulation.

\subsection{Non-Continuous Magnetic Strip}
A magnetic strip runs along the shaft of the needle, interrupted every inch
for an inch. This non-continuous configuration creates discrete magnetic
poles along the needle's length. These poles are the primary interaction
points with the solenoid array for levitation and translational control.

The choice of a non-continuous strip over a continuous one is deliberate:
discrete poles allow more precise spatial control of the magnetic interaction,
enabling finer positioning resolution in the Y (vertical) direction and
smoother transitions between solenoid activations during X and Z motion.

\subsection{No Mu-Metal Coating}
Mu-metal was considered for passive shielding of the non-strip regions of
the shaft. It was rejected on three grounds: cost, mechanical durability
(continuous knitting motion would cause wear), and physical redundancy.
Since aluminum is non-ferromagnetic, the non-strip shaft regions are already
magnetically inert without additional coating. Cross-interference from these
regions is managed through PWM frequency control and solenoid sequencing
rather than passive shielding.

%------------------------------------------------------------
\section{Solenoid Array Design}
%------------------------------------------------------------

Each actuator uses a $p \times q$ dense grid of miniature solenoids, each
individually addressable. The solenoids use silicon-steel laminated cores,
which reduce hysteresis losses and eddy current heating within the core
material itself, improving thermal efficiency.

The solenoids are driven by pulse-width modulation (PWM). PWM allows
continuous control of the effective magnetic force through duty cycle
variation without the thermal inefficiency of continuous high current.
The power supply is fully external.

Resistors and capacitors are incorporated into the driving circuit to
minimize heat loss and improve current stability.

%------------------------------------------------------------
\section{Degrees of Freedom and Control Architecture}
%------------------------------------------------------------

The needle requires full 6-DOF actuation. Three translational degrees of
freedom and three rotational degrees of freedom are identified, each
mapped to a specific control mechanism.

\subsection{Coordinate System}
A right-handed Cartesian coordinate system is defined with origin at the
center of the solenoid array surface. The Y-axis is normal to the array
surface (vertical). The X and Z axes span the horizontal plane of the array.

\subsection{DOF-to-Mechanism Mapping}

\begin{table}[h]
\centering
\begin{tabular}{>{\bfseries}l l l}
\toprule
DOF & Description & Control Mechanism \\
\midrule
X & Horizontal position (left-right) & Selective solenoid activation \\
Z & Horizontal position (forward-back) & Selective solenoid activation \\
Y & Vertical height (levitation) & PWM-controlled lifting force \\
XY plane & Tilt about Z-axis & Eddy current torque \\
YZ plane & Tilt about X-axis & Eddy current torque \\
ZX plane & Rotation about Y-axis & Eddy current torque \\
\bottomrule
\end{tabular}
\caption{Degree of freedom mapping for needle actuation.}
\end{table}

\subsection{Translational Control (X, Z)}
Horizontal needle positioning is achieved by selectively activating
solenoids ahead of and behind the needle's current position, analogous
to the operating principle of a linear motor. By controlling which
solenoids are active and at what current level, a net horizontal force
is generated that moves the needle in the XZ plane.

\subsection{Vertical Control (Y) and Earnshaw's Theorem}
Earnshaw's theorem states that a static magnetic field alone cannot
stably levitate a passive magnetic object in free space. Any equilibrium
position under purely static fields is unstable.

The system addresses this through dynamic PWM control rather than static
field equilibrium. Because the solenoid currents are continuously varied,
the field is never static. The AI camera monitoring system provides
real-time positional feedback that informs PWM adjustments, effectively
implementing active stabilization without contact sensors. Additionally,
during knitting operation the needle is never stationary --- it is
continuously in motion, which further reduces the practical impact of
Earnshaw instability compared to a static levitation scenario.

The vertical levitation condition requires that the vertical component
of the magnetic lifting force $F(t)$ exceed the gravitational load:

\begin{equation}
F'(t) > f = mg
\end{equation}

where $m$ is the mass of the needle and $g$ is gravitational acceleration.
As different solenoids activate and deactivate, $F(t)$ must be held
constant unless intentional vertical repositioning is required.

\subsection{Rotational Control (XY, YZ, ZX)}
Rotational control is achieved through eddy current torque. The
time-varying magnetic field $B(t)$ induces a magnetic flux $\phi(t)$
through the cross-sectional area $A$ of the needle's conductive region:

\begin{equation}
\phi(t) = B(t) \cdot A
\end{equation}

By Faraday's law, the changing flux induces an EMF:

\begin{equation}
e = -\frac{d\phi(t)}{dt}
\end{equation}

This drives an eddy current $E(t)$ through the aluminum core, with
effective resistance $R_t$:

\begin{equation}
E(t) \propto \frac{e(t)}{R_t}
\end{equation}

The interaction of $E(t)$ with $B(t)$ produces torque $\tau$. To achieve
a desired angular velocity $\omega$:

\begin{equation}
\tau(t) \geq I_r \frac{d\omega}{dt}
\end{equation}

where $I_r$ is the moment of inertia of the needle. Since $E(t)$ depends
directly on the rate of change of $B(t)$, the PWM drive current amplitude
and waveform are the primary control variables for rotational DOFs.

%------------------------------------------------------------
\section{Mathematical Framework}
%------------------------------------------------------------

Both actuator arrays are modeled as $p \times q$ grids of identical
miniature solenoids with silicon-steel laminated cores. Each solenoid
has length $L$, radius $r$, and $n$ turns. Each is driven by a
PWM-controlled current $c$.

The knitting needle has mass $m$ and moment of inertia $I_r$. The
array mass $M$ does not contribute to the gravitational load.

\subsection{Magnetic Force}

At any instant $t$, let $N$ solenoids participate in levitation, at
distances $r_1, r_2, \ldots, r_N$ from the needle. The average
effective distance is:

\begin{equation}
\bar{r} = \frac{1}{N} \sum_{i=1}^{N} r_i
\end{equation}

The active solenoids generate a net magnetic field $B(t)$ at the
needle's location. The needle's magnetic response generates a lifting
force $F(t)$ whose vertical component satisfies $F'(t) > mg$ whenever
levitation is required.

\subsection{Cross-Interference Between Arrays}

The complete system consists of two actuator arrays, each with its own
power supply. When both are active, each needle experiences a secondary
magnetic field $B'$ from the opposing array. If $d$ is the separation
between arrays and $\theta$ is the angle between their field vectors:

\begin{equation}
B_{\text{net}} = B + B'(\theta, d)
\end{equation}

Cross-interference is minimized through three mechanisms: PWM sequencing
(adjacent solenoids never activate in conflicting patterns), the
geometric symmetry of the two-array layout, and the natural attenuation
of field strength with distance $d$.

%------------------------------------------------------------
\section{Yarn Interaction and AI Monitoring}
%------------------------------------------------------------

Yarn properties --- tension, friction, and elasticity --- vary dynamically
during knitting and directly affect loop formation quality. Rather than
using contact-based tension sensors, this system employs a camera-based
AI monitoring approach. A vision system observes the yarn and needle
positions in real time and adjusts the solenoid array PWM parameters
accordingly.

This approach avoids the complexity of integrating physical sensors into
the work area while providing holistic feedback that captures multiple
yarn variables simultaneously. The AI system is also responsible for
active Y-axis stabilization, compensating for Earnshaw instability
through continuous field adjustment based on observed needle position.

%------------------------------------------------------------
\section{Challenges and Mitigations}
%------------------------------------------------------------

\subsection{Earnshaw Instability}
As discussed in Section 5.2, static magnetic levitation of a passive
object is inherently unstable. This is mitigated through dynamic PWM
actuation and AI-based visual feedback for active stabilization.

\subsection{Magnetic Cross-Talk}
Adjacent solenoids can interfere with each other's fields. Mitigation
strategies include: PWM sequencing to prevent conflicting adjacent
activations, array layout symmetry, and natural field attenuation with
distance. The non-ferromagnetic aluminum shaft provides passive
inertness in non-strip regions without requiring mu-metal coating.

\subsection{Thermal Management}
Continuous solenoid operation generates heat. Mitigation strategies
include: PWM duty cycle optimization (solenoids cool during inactive
periods), silicon-steel laminated cores (reduced core losses), and
increased spacing between adjacent solenoids. The external power supply
prevents heat transfer into the work area.

\subsection{Projectile Effect}
If a ferromagnetic object enters the active magnetic field, it will
experience a strong attractive force and may be launched unpredictably.
The system must be operated in a controlled environment free of
ferromagnetic materials.

\subsection{Strong Magnetic Field}
Prolonged human exposure to the operating magnetic field should be
avoided. Appropriate shielding of the external environment is recommended.

%------------------------------------------------------------
\section{Questions and Answers}
%------------------------------------------------------------

\textbf{Won't gravity restrict the needle from levitating?}

The gravitational force on the needle does not prevent levitation. The
vertical component of the magnetic lifting force is calibrated to
satisfy $F'(t) > mg$, directly counteracting the needle's weight.

\textbf{What happens to the knitted fabric?}

The knitted fabric is collected in a wooden tray positioned below the
work area. The tray height is determined by the breaking stress of the
needle and the linear mass density $\alpha$ of the chosen knitting
pattern. If $W$ is the maximum weight the needle can support without
bending, the tray is placed at distance $W / \alpha$ below the needle's
lowest operating position.

%------------------------------------------------------------
\section{Limitations}
%------------------------------------------------------------

\begin{itemize}
    \item The mathematical framework presented here is theoretical.
    No prototype has been built and no empirical validation has been
    conducted.
    \item PWM frequency optimization for simultaneous torque generation
    and cross-interference suppression requires detailed simulation and
    experimental tuning.
    \item Yarn physics, including the dynamic interaction between yarn
    and needle during loop formation, requires further characterization.
    \item Active noise cancellation may be required, as rapid solenoid
    switching generates acoustic noise.
    \item A backup power supply is necessary: sudden power interruption
    can cause needle drop and potential solenoid damage.
\end{itemize}

%------------------------------------------------------------
\section{Future Work}
%------------------------------------------------------------

\begin{itemize}
    \item \textbf{Prototype development}: Construction and empirical
    validation of the two-needle system with a representative solenoid
    array.
    \item \textbf{Numerical simulation}: Finite element analysis of
    the magnetic field distribution, eddy current torque, and
    cross-interference for the proposed needle geometry.
    \item \textbf{AI monitoring system}: Development and training of
    the camera-based yarn monitoring and active stabilization system.
    \item \textbf{Multi-needle extension}: Scaling the architecture
    to a multi-needle system for continuous knitting machine operation.
    \item \textbf{Pattern complexity}: Integration of AI-driven pattern
    generation for complex knit structures.
\end{itemize}

%------------------------------------------------------------
\section*{References}
%------------------------------------------------------------

\begingroup
\renewcommand{\section}[2]{}
\begin{thebibliography}{9}

\bibitem{zhang2019}
Zhang, C., Li, D., Zuo, X., \& Zhou, X. (2019).
Magnetic array knitting needle driving system design.
\textit{IOP Conference Series: Materials Science and Engineering},
612(3), 032026.

\bibitem{zuo2021}
Zuo, X., Zhang, C., Xiong, T., Yin, W., Li, M., Zhang, C., Wu, X.,
\& Zhu, L. (2021).
Maglev weft knitting needle driving modeling and PID controller design.
\textit{The Journal of The Textile Institute}, 113(8), 1715--1722.

\bibitem{xiong2023}
Xiong, T., Peng, Y., Zuo, X., Zhang, C., Zhang, C., Zhang, L.,
\& Li, H. (2023).
Magnetic field analysis and optimization of the gauge of hybrid maglev
needles.
\textit{Applied Sciences}, 13(3), 1257.

\bibitem{zhang2025}
Zhang, D., Sheng, X., Gao, P., Wang, C., \& Shi, Y. (2025).
Structure design and control of magnetic levitated needle driving system.
In \textit{Intelligent Robotics and Applications, ICIRA 2024},
Lecture Notes in Computer Science, vol.\ 15207. Springer, Singapore.

\end{thebibliography}
\endgroup

\end{document}
